\section{Recipe 2：Zephyr 与 Alignment Distillation}

Human data 昂贵，课堂随后转向第二条 recipe：让强模型产生 dialogue 与 preference，再把 alignment behavior 蒸馏到较小 open model。这个路线可以扩展，但 teacher model 的风格与 judge 偏差也会被复制。

\subsection{从 Human Feedback 到 AI Feedback}

\term{alignment distillation（对齐蒸馏）}在这里指用 AI-generated dialogues 与 AI-ranked preferences 训练较小模型。它不是把 teacher logits 直接做传统 knowledge distillation，而是蒸馏 instruction-following 与 preference behavior。

\lecturefigure{slide-30.jpg}{Recipe 2：Distillation of AI Alignment}{官方幻灯片 p.30}

\begin{importantbox}{两条 recipe 的差异在 label producer}
Human recipe 的 demonstrations/preferences 主要由人生产；Zephyr recipe 的 dialogues 与 rankings 主要由强模型生产。两者可使用相似 SFT/DPO objective，但监督信号的 provenance 和 failure correlation 完全不同。
\end{importantbox}

\lecturefigure{slide-31.jpg}{Zephyr-7B：dSFT、AI feedback 与 dDPO 三阶段}{官方幻灯片 p.31}

\begin{knowledgebox}{读图：三列是一条数据转换管线}
第一列从 prompt 生成 multi-turn dialogues，用于 distilled SFT；第二列让多个模型生成候选并由 GPT-4 ranking；第三列把 chosen/rejected pair 送入 DPO。每一列都会筛选分布，也会继承 teacher/judge 偏差。
\end{knowledgebox}

\term{DPO（Direct Preference Optimization）}直接比较 policy 对 chosen/rejected 的相对 log-probability，并用 reference policy 校准：
\[
\mathcal{L}_{\mathrm{DPO}}(\theta)
=-\log\sigma\!\left(
\beta\log\frac{\pi_{\theta}(y_w\mid x)}{\pi_{\mathrm{ref}}(y_w\mid x)}
-\beta\log\frac{\pi_{\theta}(y_l\mid x)}{\pi_{\mathrm{ref}}(y_l\mid x)}
\right).
\]
其中 $\beta$ 控制偏好更新强度，$y_w/y_l$ 分别是 chosen/rejected。DPO 不显式训练 reward model，也不等于“无需 SFT”；后面的 ablation 恰恰显示 SFT 是关键起点。

\teachervoice{Rajani 把 Zephyr 称为“distillation of AI alignment”：UltraChat 负责 dSFT，AI-generated rankings 形成 UltraFeedback，最后用 DPO 蒸馏 preferences。课堂结果必须同时记住 teacher/judge 都来自强模型。}

\subsection{本章小结}

Zephyr 用 AI data 替代大量人工写作与排序，把 alignment pipeline 变成可扩展的 distillation。成本下降的同时，training--evaluation coupling 和 teacher style imitation 成为新的主要风险。

